\documentclass[10pt,a4paper,landscape]{article}
\usepackage[margin=0.5cm]{geometry}
\usepackage{multicol}
\usepackage{amsmath,amssymb}
\usepackage{enumitem}
\usepackage{tikz}
\usepackage{flowchart}
\usepackage[utf8]{inputenc}
\usepackage[japanese]{babel}
\usepackage{xeCJK}

\setlength{\columnseprule}{0.4pt}
\setlist{nosep}
\pagestyle{empty}

\begin{document}

% ページ1
\begin{multicols}{2}

\section*{基本概念と線形代数}

\subsection*{最適化問題の基本形}
\begin{align}
\min/\max \quad & f(x) \\
\text{s.t.} \quad & g_i(x) \leq b_i \quad \text{(不等式制約)} \\
& h_j(x) = c_j \quad \text{(等式制約)} \\
& x \in X \quad \text{(変数の範囲)}
\end{align}

\subsection*{線形代数の基礎}
\begin{itemize}
\item \textbf{階数 (Rank)}: 独立な行/列の最大数
\item \textbf{行列積}: $C = AB$, $c_{ij} = \sum_k a_{ik}b_{kj}$
\item \textbf{逆行列}: $AA^{-1} = I$
\item \textbf{連立方程式} $Ax = b$:
  \begin{itemize}
  \item $m = n$, $\text{rank}(A) = n \rightarrow$ 一意解
  \item $m < n \rightarrow$ 無限解
  \item $\text{rank}(A) < \text{rank}([A|b]) \rightarrow$ 解なし
  \end{itemize}
\end{itemize}

\subsection*{勾配とヘシアン}
\begin{itemize}
\item \textbf{勾配}: $\nabla f(x) = \begin{bmatrix} \frac{\partial f}{\partial x_1} \\ \vdots \\ \frac{\partial f}{\partial x_n} \end{bmatrix}$
\item \textbf{ヘシアン}: $\nabla^2 f(x) = \left[\frac{\partial^2 f}{\partial x_i \partial x_j}\right]$
\item \textbf{テイラー1次}: $f(x) \approx f(x_0) + \nabla f(x_0)^T(x - x_0)$
\item \textbf{テイラー2次}: $f(x) \approx f(x_0) + \nabla f(x_0)^T(x - x_0) + \frac{1}{2}(x - x_0)^T\nabla^2 f(x_0)(x - x_0)$
\end{itemize}

\subsection*{最適性条件}
\begin{itemize}
\item \textbf{1次必要条件}: $\nabla f(x^*) = 0$
\item \textbf{2次十分条件}: 
  \begin{itemize}
  \item min問題: $\nabla^2 f(x^*)$ が正定値 (PD)
  \item max問題: $\nabla^2 f(x^*)$ が負定値 (ND)
  \end{itemize}
\end{itemize}

\columnbreak

\section*{凸性判定フローチャート}

\subsection*{凸集合 (Convex Set) 判定}
\begin{center}
\begin{tikzpicture}[node distance=1.5cm]
\node (start) [rectangle, draw] {集合 $S$ は凸か？};
\node (test) [rectangle, draw, below of=start] {$\forall x,y \in S, \forall \lambda \in [0,1]:$ \\ $\lambda x + (1-\lambda)y \in S$ か？};
\node (yes) [rectangle, draw, below left of=test] {YES $\rightarrow$ 凸集合};
\node (no) [rectangle, draw, below right of=test] {NO $\rightarrow$ 非凸集合};

\draw [arrow] (start) -- (test);
\draw [arrow] (test) -- (yes);
\draw [arrow] (test) -- (no);
\end{tikzpicture}
\end{center}

\textbf{凸集合の例}:
\begin{itemize}
\item 超平面: $\{x | a^T x = b\}$
\item 半空間: $\{x | a^T x \leq b\}$
\item ボール: $\{x | \|x - c\| \leq r\}$
\item 多面体: $Ax \leq b$
\end{itemize}

\subsection*{凸関数 (Convex Function) 判定}
\begin{center}
\begin{tikzpicture}[node distance=1.2cm]
\node (start) [rectangle, draw] {関数 $f(x)$ は凸か？};
\node (def) [rectangle, draw, below of=start, text width=4cm] {【定義】$\forall x,y, \forall \lambda \in [0,1]:$ \\ $f(\lambda x + (1-\lambda)y) \leq \lambda f(x) + (1-\lambda)f(y)$};
\node (first) [rectangle, draw, below of=def, text width=4cm] {【1次条件】$\forall x,y:$ \\ $f(y) \geq f(x) + \nabla f(x)^T(y - x)$};
\node (second) [rectangle, draw, below of=first, text width=3cm] {【2次条件】\\ $\nabla^2 f(x)$ が正半定値};
\node (result) [rectangle, draw, below of=second] {いずれかYES $\rightarrow$ 凸関数};

\draw [arrow] (start) -- (def);
\draw [arrow] (def) -- (first);
\draw [arrow] (first) -- (second);
\draw [arrow] (second) -- (result);
\end{tikzpicture}
\end{center}

\textbf{正半定値の判定}:
\begin{itemize}
\item 全固有値 $\geq 0$
\item 主小行列式が非負
\item $x^T A x \geq 0 \, \forall x$
\end{itemize}

\end{multicols}

\newpage

% ページ2
\begin{multicols}{2}

\section*{凸計画問題と線形計画法}

\subsection*{凸計画問題判定}
\begin{center}
\begin{tikzpicture}[node distance=1.3cm]
\node (start) [rectangle, draw, text width=3.5cm] {最適化問題は凸計画問題か？};
\node (min) [rectangle, draw, below of=start, text width=4cm] {【min問題】\\ 1. 目的関数 $f(x)$ が凸関数\\ 2. 実行可能領域が凸集合};
\node (max) [rectangle, draw, right of=min, text width=4cm] {【max問題】\\ 1. 目的関数 $f(x)$ が凹関数\\ 2. 実行可能領域が凸集合};
\node (result1) [rectangle, draw, below of=min] {両方YES $\rightarrow$ 凸計画問題};
\node (result2) [rectangle, draw, below of=max] {いずれかNO $\rightarrow$ 非凸};

\draw [arrow] (start) -- (min);
\draw [arrow] (start) -- (max);
\draw [arrow] (min) -- (result1);
\draw [arrow] (max) -- (result2);
\end{tikzpicture}
\end{center}

\subsection*{線形計画法 (Linear Programming)}
\textbf{標準形}:
\begin{align}
\min \quad & c^T x \\
\text{s.t.} \quad & Ax = b \\
& x \geq 0
\end{align}

\textbf{基本用語}:
\begin{itemize}
\item \textbf{基底解}: $n$個の基底変数の解
\item \textbf{基本実行可能解 (BFS)}: 基底解かつ実行可能
\item \textbf{退化}: 基本変数の値が0
\item \textbf{極端点}: 多面体の頂点
\item \textbf{極端線}: 無限に延びる方向
\end{itemize}

\subsection*{シンプレックス法}
\begin{enumerate}
\item \textbf{初期BFS}を見つける
\item \textbf{最適性判定}: 被約費用 $\leq 0$ なら最適
\item \textbf{方向決定}: 入る変数を選択
\item \textbf{ステップサイズ}: 最小比率テスト
\item \textbf{基底更新}: 出る変数を決定
\end{enumerate}

\columnbreak

\section*{特殊問題と実用公式}

\subsection*{輸送問題 (Transportation Problem)}
\begin{align}
\min \quad & \sum_i \sum_j c_{ij} x_{ij} \\
\text{s.t.} \quad & \sum_j x_{ij} \leq s_i \quad \text{(供給制約)} \\
& \sum_i x_{ij} \geq d_j \quad \text{(需要制約)} \\
& x_{ij} \geq 0
\end{align}

\textbf{実行不能条件}: $\sum s_i < \max\{d_j\}$

\subsection*{絶対値の線形化}
\begin{align}
\min |x - a| \quad \rightarrow \quad & \min z \\
\text{s.t.} \quad & z \geq x - a \\
& z \geq -(x - a)
\end{align}

\subsection*{区間制約の判定}
集合 $X = \{x \in [a,b]^n : \sum x_i \geq c\}$
\begin{itemize}
\item \textbf{有界}: $[a,b]$が有限区間
\item \textbf{閉}: 等式・不等式制約で定義
\item \textbf{凸}: 箱制約と線形制約の交集合
\end{itemize}

\subsection*{よく出る凸関数}
\begin{itemize}
\item $x^k$ ($k \geq 1$, $x \geq 0$): 凸
\item $\ln(x)$ ($x > 0$): 凹
\item $e^x$: 凸
\item $x^2 + y^2$: 凸
\item $\max\{f_1(x), f_2(x)\}$: $f_1, f_2$が凸なら凸
\item $|x|$: 凸
\end{itemize}

\subsection*{決定変数の制約}
\begin{itemize}
\item \textbf{"here-and-now"}: 不確実性実現前に決定
\item \textbf{"wait-and-see"}: 不確実性実現後に決定
\end{itemize}

\subsection*{重要な数値例}
過去問から:
\begin{itemize}
\item \textbf{勾配計算}: $f(x_1,x_2) = \sqrt{x_1-1} + (x_2-1)^{1/4}$ at $(3,2)$
  \[\nabla f = \begin{bmatrix} \frac{1}{2\sqrt{2}} \\ \frac{1}{4} \end{bmatrix}\]
\item \textbf{ヘシアン計算}: 各偏導関数の2次偏微分
\end{itemize}

\end{multicols}

\vfill

\section*{重要記号・略語}
\begin{itemize}
\item \textbf{LP}: Linear Programming (線形計画法)
\item \textbf{NLP}: Nonlinear Programming (非線形計画法)  
\item \textbf{BFS}: Basic Feasible Solution (基本実行可能解)
\item \textbf{PSD}: Positive Semidefinite (正半定値)
\item \textbf{PD}: Positive Definite (正定値)
\item \textbf{s.t.}: subject to (制約条件)
\end{itemize}

\end{document}